%%FILE I-a
%%APPEND title="Support, zero extension, and the exceptional inverse image"
In a pointed derivator the two chains of adjoints of this page exist without the stability axiom: for an open inclusion $j$ (a sieve) and the complementary closed inclusion $i$ (a cosieve) one has $j^?\dashv j_!\dashv j^*\dashv j_*$ and $i_!\dashv i^*\dashv i_*\dashv i^!$, the four Kan extensions are fully faithful, and their images are described by supports; the coherent square relating the two extensions by zero is constructed by Kan extensions on the product with a square. This is in \citet{Groth2013}; the numbers of the statements there were not checked for this note. Stability is needed only to turn the resulting fibre and cofibre sequences into the triangles of a recollement.

%%BLOCK after="From derived coefficients to diagram categories"
\begin{dcomment}{status}{p4}{Ordered sets as stratified spaces: what is proved}
The topological vocabulary of these pages (open and closed immersions, proper and smooth maps, direct and exceptional images) is used for functors between small categories. For partially ordered sets part of it is a theorem about spaces, and part of it is not; the precise state is as follows. Give a poset $P$ the topology whose open sets are the downward closed subsets, so that a presheaf on $P$ in the sense of the manuscript has the values of a sheaf on the principal open sets.

(1) For every poset $P$, presheaves of spaces on $P$ are the hypercomplete sheaves of spaces on this topological space \citep[Proposition~3.12.11 and Example~3.12.15]{BGH2020}; they are all the sheaves when the principal open sets are finite. Without a finiteness hypothesis the sheaves form a different category: for the opposite of the ordinal sum of countably many two-point antichains the terminal sheaf is a compact object, and the terminal presheaf is not \citep[Example~A.13]{Aoki2023}. For presheaves of sets there is no restriction.

(2) For a finite poset $P$, the classifying space $|NP|$, stratified over $P$ by the least vertex of the open simplex containing a point, has contractible strata, and the sheaves of spaces on $|NP|$ that are locally constant on the strata are again the presheaves of spaces on $P$. The identification passes through a comparison of two stratifications of $|NP|$ whose source was not opened for this commentary; it is recorded as stated.

(3) A derivator supplies the categories $\mathbb D(P)$ and their Kan extensions, and no sheaf theory on a space; (1) and (2) concern the derivator of spaces, and extensions to other coefficients need hypotheses on the coefficients.

(4) The proper and smooth functors of ms.~52--54 are conditions on categories of lifts. For a monotone map of finite posets they are neither the proper maps nor the submersions of the realizations: the two inclusions of an endpoint in $[1]$ are homeomorphic closed embeddings of realizations, and one is proper and not smooth, the other smooth and not proper. A full inclusion of posets is proper exactly when it is a cosieve and smooth exactly when it is a sieve; product projections are both.

(5) The Kan extensions along a monotone map of finite posets are the topological direct images of the corresponding sheaves for closed inclusions and for product projections. For an arbitrary monotone map they are not; whether they are for all proper maps in the sense of (4) has been asserted and was not verified for this commentary.

(6) A category in which every arrow is a monomorphism, every isomorphism an identity and every endomorphism an identity, with finitely many arrows, still has a realization stratified over the poset of its objects, in which parallel arrows give distinct cells (with the same unopened comparison as in (2)); the category with three objects and nine arrows whose classifying space is the real projective plane with its flag stratification is an example. For a base with a noninvertible endomorphism, such as the monoid $\{1,e\}$ with $e^2=e$, the presheaves of spaces are not the constructible sheaves of any stratified space over a poset: the idempotent survives in every equivalent presheaf category, although the classifying space is contractible.
\end{dcomment}

%%APPEND title="Properness as a condition on all lifts"
For coefficients in spaces the base-change property characterizes properness. If $f:X\to Y$ satisfies right base change along every strict pullback, apply it to the terminal coefficient: for an object $x$ and an arrow $f(x)\to y$ the comparison map is the map from the disjoint union of the nerves of the categories of lifts to the discrete set $\operatorname{Hom}_Y(f(x),y)$, and it is an equivalence exactly when every category of lifts is weakly contractible.

%%FILE VI-a
%%APPEND title="Path completion and the modern category of local systems"
How much of the base the two theories remember can be said exactly when the coefficients are enhanced. The category of all space-valued diagrams on $S$, with its enhancement, determines the idempotent completion of $S$ and nothing more: the point and the interval have equivalent classifying spaces and inequivalent diagram categories (two and three subterminal objects). The locally constant theory together with the constant-diagram functor determines the homotopy type $BS$. Without the enhancement or the constant-diagram functor the reconstruction is not known. On the interval $[1]$ the reflection and the coreflection onto locally constant objects take the values of a diagram at the two opposite endpoints.

%%FILE VI-b
%%APPEND title="The extra structure hidden in a categorical base"
The two fine localizations over a fixed base are not equivalent to each other in general: they are the theories of covariant and of contravariant diagrams of spaces on $S$, exchanged by passing to $S^{\mathrm{op}}$, and for the fork (one object with arrows to two others) the category is idempotent complete and not equivalent to its opposite, so the two diagram theories differ. Exchanging them is a change of base, not a duality on a fixed base. For a base with a noninvertible endomorphism the fine theory is moreover not the constructible theory of any space stratified over a poset (commentary on Chapter~I, ms.~4).

%%FILE XIIb-a
%%APPEND title="The presheaf case is already present in SGA 4"
With coefficients in an $\infty$-category $\mathcal C$ the same gluing holds coherently. For a cut of $X$ into a sieve $A$ and the complementary cosieve $B$, with $E$ the category of crossing arrows and $s:E\to A$, $t:E\to B$ its two projections, the gluing functor $i^*j_*$ is $t_*s^*$, and $\operatorname{Fun}(X^{\mathrm{op}},\mathcal C)$ is the lax pullback of triples $(F_A,F_B,F_B\to t_*s^*F_A)$; the locally constant objects form the pullback $\operatorname{Fun}(BA,\mathcal C)\times_{\operatorname{Fun}(BE,\mathcal C)}\operatorname{Fun}(BB,\mathcal C)$. The construction must be made before passing to homotopy categories: already over $[1]$ the comma category formed from homotopy categories has one morphism from $S^0\to*$ to $*\to S^1$ where the coherent arrow category has a set of classes indexed by $\mathbf Z$. When the inclusion is discretely flat in the sense used in the commentary at ms.~114, the strict pushout of categories is also the pushout of the associated $\infty$-categories.

%%FILE XIIapp
%%APPEND title="The open--closed gluing formula"
\textbf{Complements.}
(1) The category $E$ of crossing arrows has a description in the stratified space $BX$: the simplicial set of the maps $\Delta^1\to N(X^{\mathrm{op}})$ over $\Delta^1$ is the nerve of $E^{\mathrm{op}}$, and $BE$ is weakly equivalent, compatibly with the two end-point maps, to the homotopy link of $BB$ in $BX$, the space of paths that start in $BB$ and leave it immediately \citep[Theorem~6.1]{CDDouteauWaas2022}; the complement of $BB$ in $BX$ deformation retracts onto $BA$. So the gluing formula is the decomposition of a two-strata space along its homotopy link.
(2) For a finite filtration of $X$ by sieves $A_0\subset\dots\subset A_n=X$ the formula applies to each step $A_{k-1}\subset A_k$; the homotopy types of the pairwise crossing categories do not suffice to recover $BX$ (the boundary of a triangle and the filled triangle have the same ones), and the categories of flags of crossing arrows are needed.
(3) The formula can also be proved at the level of categories, without nerves: let $P$ be the category over the three-object category $\cdot\leftarrow\cdot\to\cdot$ with fibres $A$, $E$, $B$; for the evident functor $P\to X$ and an object $x$ of $B$ the category $P/x$ carries an endofunctor with natural transformations from the identity and from the constant functor at $(x,\mathrm{id}_x)$, and for $x$ in $A$ it is $A/x$. Theorem~A gives the result for every fundamental localizer satisfying the relative axiom; no reverse excision is used.

%%FILE XIIc-a
%%APPEND title="Degree two and the localizer correction"
The intrinsic fibrations of the localizer $W_n$ of the maps satisfying $b_n$ are known: a functor $p:T\to B$ preserves $W_n$ under every base change if and only if every strict pullback square along $p$ becomes a homotopy pullback of $n$-truncated types \citep[Proposition~6.2.3 and \S9.2]{Cisinski2006}. A fibration that is also an opfibration need not have this property: the functor from the translation groupoid of the group with two elements to that group is one, the functor $e\to BC_2$ is in $W_0$, and its pullback is the inclusion of a two-point discrete category in a contractible one.

%%FILE XVI
%%APPEND title="Why enough intrinsic fibrations gives B"
The factorizations of this section produce functors that preserve weak equivalences under base change, and perfect or locally constant weak fibrations; they do not produce functors that are fibrations and opfibrations in the strict sense. That question is the one of Chapter~VI, ms.~186.

%%FILE XI-XIIa
%%APPEND title="From strict fibres to homotopy fibres"
A $\mathsf{Hot}$-fibration has a proper and smooth replacement with the same fibres up to weak equivalence, by the alternating construction of Chapter~VI, ms.~192--195. A replacement that is a Grothendieck fibration and opfibration is known when the base is a poset, or more generally has only monomorphisms or only epimorphisms, when its classifying space is a 1-type, and when the fibres are 1-types (commentary on Chapter~VI, ms.~186); in general it is not known.
